\section{Conclusión}
% Que puedo decir con base en lo que obtuve

Con base en lo mostrado en las secciones antriores podemos concluir que el algoritmo de colonia de hormigas (ACO) es una
herramienta efectiva para abordar el problema de bin packing. Los resultados indican que el ACO puede encontrar soluciones cercanas al óptimo en un tiempo razonable, 
demostrando su capacidad para explorar y explotar el espacio de soluciones de manera eficiente.

Si tenemos en cuenta el maximo de contenedores posibles (401) y el numero de contenedores usados (161) podemos ver que tenemos una reduccion del 59.85\% es decir un poco más
de la mitad de los contenedores posibles, lo cual es un resultado bastante positivo. Si traducimos esto en terminos de costos logísticos, el uso eficiente de los contenedores
puede resultar en ahorros significativos, ya que se reduce la cantidad de espacio desperdiciado y se optimizan los recursos disponibles. Ademas de esto como pudimos ver en las
figuras de la sección de discusión, la distribución de cargas en los contenedores y columnas es bastante uniforme, sin embargo todavía hay espacio para mejorar la eficiencia 
del empaquetado lo que a su vez podría llevar a una reducción adicional en el número de contenedores utilizados.

Tambien podemos ver que las distribución es bastante uniforme, lo que sugiere que el ACO es capaz de distribuir los paquetes de manera efectiva entre los contenedores aun cuando
el tipo de paquetes es aleatorio. Esto es importante ya que en escenarios del mundo real, la demanda de paquetes puede variar significativamente, y un algoritmo robusto debe ser capaz
de manejar esta variabilidad.
